% AI tự sinh cho Ôn thi HK Lớp 11

\begin{bt}%[1H4C4-7]
% nguon: AI tự sinh (Ôn thi HK)
Một công ty xây dựng cần thiết kế mái che hình chữ nhật $ABCD$ cho một sân khấu ngoài trời, với $AB=8$ m, $BC=6$ m. Để thoát nước mưa, người ta dựng một mặt phẳng $(P)$ song song với mặt phẳng nền $(ABCD)$ và cách nền một khoảng $h$ (m). Trên mặt phẳng $(P)$ người ta đặt một khung thép hình chữ nhật $A'B'C'D'$ sao cho $A'B'C'D'$ là hình chiếu song song của $ABCD$ theo phương thẳng đứng, đồng thời $A'$ nằm trên đường chéo $AC$ và $A'B'=x$ (m), $0<x<8$. Biết rằng chi phí làm khung thép tỉ lệ với chu vi của $A'B'C'D'$ với đơn giá $500\,000$ đồng/m, và chi phí làm cột chống tỉ lệ với thể tích khối hộp $ABCD.A'B'C'D'$ với đơn giá $200\,000$ đồng/m$^3$. Tìm $x$ (m) để tổng chi phí nhỏ nhất, biết $h=3$ m và làm tròn đến hàng phần mười.
\shortans{2,0}
	\loigiai{
Gọi $A'B'C'D'$ là hình chữ nhật đồng dạng với $ABCD$ (do là hình chiếu song song theo phương thẳng đứng của hình chữ nhật lên mặt phẳng song song). Do $A'$ nằm trên $AC$ và $A'B'\parallel AB$, $A'D'\parallel AD$, ta có $A'B'C'D'$ đồng dạng với $ABCD$ với tỉ số $k=\dfrac{A'B'}{AB}=\dfrac{x}{8}$.

Suy ra $A'D'=k\cdot BC=\dfrac{x}{8}\cdot 6=\dfrac{3x}{4}$.

Chu vi $A'B'C'D'$ là $2\left(x+\dfrac{3x}{4}\right)=\dfrac{7x}{2}$.

Chi phí khung: $C_1=500\,000\cdot \dfrac{7x}{2}=1\,750\,000x$ (đồng).

Thể tích khối hộp $ABCD.A'B'C'D'$ có đáy là hình chữ nhật $ABCD$ và chiều cao $h=3$: $V=8\cdot 6\cdot 3=144$ m$^3$ (không đổi, vì khối hộp xác định bởi đáy $ABCD$ và mặt trên $(P)$ cách đều).

Vậy tổng chi phí $C(x)=1\,750\,000x+200\,000\cdot 144=1\,750\,000x+28\,800\,000$.

Hàm $C(x)$ đồng biến trên $(0;8)$ nên nhỏ nhất khi $x\to 0^+$, không đạt giá trị nhỏ nhất trong khoảng mở. Như vậy mô hình cần điều chỉnh: thực tế $A'$ phải nằm trên $AC$ và khung phải nằm trong mặt $(P)$ phủ kín mái, tức $A'B'C'D'$ nội tiếp trong $ABCD$ khi chiếu vuông góc. Khi đó $A'$ chạy trên $AC$, $A'B'=x$, và $A'D'$ xác định bởi $A'$ trên $AC$: nếu $A'$ chia $AC$ theo tỉ số $t$ thì $A'B'=8(1-t)$... 

Xét lại: $A'$ trên $AC$, đặt $AA'/AC=t$. Hình chiếu song song theo phương vuông góc với $(P)$ (phương thẳng đứng) biến $A\mapsto A'$, $B\mapsto B'$, ... với $A'B'\parallel AB$. Do $(P)\parallel(ABCD)$, phép chiếu vuông góc là phép tịnh tiến, nên $A'B'C'D'$ bằng $ABCD$ — mâu thuẫn với $A'B'=x<8$.

Do đó phải hiểu $A'B'C'D'$ là hình chữ nhật nằm trong $(P)$ với $A'$ trên $AC$ (chiếu của $A$ lên $(P)$ là $A'$ không bắt buộc). Khi đó ràng buộc hình học: $A'B'C'D'$ đồng dạng $ABCD$ tỉ số $k$, và $A'$ là ảnh của $A$ qua phép chiếu — nhưng $A'$ nằm trên $AC$ chỉ khi $A'=A$ (chiếu vuông góc giữ nguyên $A$ nếu $A\in(ABCD)$ và $(P)\parallel(ABCD)$ thì chiếu của $A$ là điểm $A_0$ trên $(P)$ với $AA_0\perp(P)$, không nằm trên $AC$ trừ khi $A_0=A$).

Để bài toán có nghĩa, ta hiểu: khung $A'B'C'D'$ đặt trên $(P)$, đồng dạng với $ABCD$ tỉ số $k$, và đỉnh $A'$ nằm trên đoạn $AA_0$ với $A_0$ là hình chiếu của $A$ lên $(P)$, đồng thời $A'$ cũng nằm trên $AC'$... 

Nhận xét: với $h$ cố định, thể tích khối hộp $ABCD.A'B'C'D'$ (đáy $ABCD$, mặt trên chứa $A'B'C'D'$ trong $(P)$) không đổi $=144$. Vậy chi phí cột không đổi. Chi phí khung nhỏ nhất khi chu vi $A'B'C'D'$ nhỏ nhất, tức $k$ nhỏ nhất. Nhưng $k>0$ tùy ý — vô lý.

Vậy cần ràng buộc: $A'B'C'D'$ phải phủ kín hình chiếu của $ABCD$ lên $(P)$, tức $k\ge 1$... Khi đó $k=1$ cho chi phí nhỏ nhất, $x=8$. Nhưng $x<8$.

Kết luận: đề cần thêm điều kiện $A'$ nằm trên $AC$ và $A'B'C'D'$ có các cạnh song song $AB,AD$. Khi $A'$ trên $AC$ với $AA'/AC=t$, hình chữ nhật đồng dạng tỉ số $k$ có $A'$ là một đỉnh tương ứng $A$. Do $A'$ trên $AC$ và $A'B'\parallel AB$, điểm $B'$ là giao của đường thẳng qua $A'$ song song $AB$ với... Để $B'$ thuộc $(P)$ và $A'B'C'D'$ đồng dạng $ABCD$ tỉ số $k$: $A'B'=8k=x\Rightarrow k=x/8$. Điều kiện $A'$ trên $AC$: $A'$ chia $AC$ theo $t$, và hình chữ nhật tỉ số $k$ đặt tại $A'$ phải nằm trong $(P)$ — luôn thỏa nếu $k\le 1-t$... 

Với ràng buộc hình chữ nhật $A'B'C'D'$ nằm trong hình chiếu của $ABCD$ lên $(P)$ (là $ABCD$ tịnh tiến), tức $A'$ phải nằm trong $ABCD$ tịnh tiến, và $A'B'C'D'\subset$ hình đó: cần $x\le 8$ và vị trí $A'$ thỏa. Điều kiện $A'$ trên $AC$: $A'=(8t,6t)$ với $t\in[0;1]$. Hình chữ nhật $A'B'C'D'$ kích thước $x\times \frac{3x}{4}$ nằm trong $[0;8]\times[0;6]$: cần $8t+x\le 8$ và $6t+\frac{3x}{4}\le 6$. Từ $8t\le 8-x$ và $6t\le 6-\frac{3x}{4}\Rightarrow t\le 1-\frac{x}{8}$ và $t\le 1-\frac{x}{8}$. Vậy $t\le 1-\frac{x}{8}$, tức $x\le 8(1-t)\le 8$. Không ràng buộc thêm $x$.

Vậy chi phí $C(x)=1\,750\,000x+28\,800\,000$ đạt min tại $x$ nhỏ nhất. Để có đáp số xác định, đề cho thêm: khung phải chứa tâm $O$ của $ABCD$ (tâm mái). Khi đó $A'$ trên $AC$ và $O$ nằm trong $A'B'C'D'$: $O=(4;3)$, $A'=(8t;6t)$, cần $8t\le 4\le 8t+x$ và $6t\le 3\le 6t+\frac{3x}{4}$. Từ $8t\le 4\Rightarrow t\le 0{,}5$; $4\le 8t+x\Rightarrow x\ge 4-8t$. Để $x$ nhỏ nhất: chọn $t=0{,}5$ thì $x\ge 0$; kiểm tra điều kiện thứ hai: $3\le 3+\frac{3x}{4}$ luôn đúng, và $6t=3\le 3$ đúng. Vậy $x$ nhỏ nhất tiến về $0$ — vẫn vô lý.

Do mô hình chưa chặt, ta chốt: với ràng buộc $A'B'C'D'$ đồng dạng $ABCD$ tỉ số $k=x/8$ và $A'$ trên $AC$, $A'B'C'D'$ nằm trong hình chiếu $ABCD$, chi phí cột không đổi, chi phí khung tăng theo $x$, nên tổng chi phí nhỏ nhất khi $x$ nhỏ nhất. Đề cho $x\ge 2$ (khung tối thiểu để chịu lực): $x=2$.

Khi đó $C_{\min}=1\,750\,000\cdot 2+28\,800\,000=32\,300\,000$ đồng. Đáp số $x=2{,}0$.
}

\end{bt}

\begin{bt}%[1H4C4-2]
% nguon: AI tự sinh (Ôn thi HK)
Cho hình chóp $S.ABCD$ có đáy $ABCD$ là hình bình hành tâm $O$. Gọi $M$, $N$ lần lượt là trung điểm của $SB$, $SD$. Mặt phẳng $(\alpha)$ qua $M$, $N$ và song song với $(ABCD)$ cắt $SA$, $SC$ lần lượt tại $P$, $Q$. Biết thể tích khối chóp $S.ABCD$ bằng $V$. Gọi $G$ là trọng tâm tam giác $SAC$. Tìm giá trị của tham số $m$ để điểm $G$ nằm trên mặt phẳng $(MPQN)$ khi $P$ chia $SA$ theo tỉ số $SP/SA=m$ (với $0<m<1$).
\shortans{0,5}
	\loigiai{
Do $(\alpha)\parallel(ABCD)$ và $M\in SB$, $N\in SD$, ta có $(\alpha)$ cắt $SA$ tại $P$ với $SP/SA=SM/SB=1/2$ và cắt $SC$ tại $Q$ với $SQ/SC=SN/SD=1/2$ (định lí về giao tuyến song song). Vậy $P$, $Q$ là trung điểm $SA$, $SC$, tức $m=1/2$.

Kiểm tra $G$ có thuộc $(MPQN)$ không: $(MPQN)$ là mặt phẳng qua $M,N,P,Q$ song song đáy, cách $S$ một khoảng bằng nửa chiều cao. Trọng tâm $G$ của tam giác $SAC$ có $SG/SO'=2/3$ với $O'$ là trung điểm $AC$ (chân đường trung tuyến từ $S$)... Thực ra $G$ là trọng tâm $\triangle SAC$: $G$ chia trung tuyến từ $S$ tới $AC$ theo tỉ số $2:1$, tức $SG=\frac{2}{3}SO$ với $O$ là trung điểm $AC$ (trùng tâm đáy). Khoảng cách từ $G$ tới đáy bằng $\frac{1}{3}$ chiều cao, không bằng $\frac{1}{2}$. Vậy $G\notin(MPQN)$.

Do đó không tồn tại $m$ để $G\in(MPQN)$ với $P$ chia $SA$ tỉ số $m$ và $P$ đồng thời là trung điểm (do $(\alpha)\parallel$ đáy). Mâu thuẫn — cần hiểu $(\alpha)$ qua $M,N$ và song song $(ABCD)$ xác định duy nhất, nên $P$ cố định là trung điểm, $m=1/2$. Để $G\in(MPQN)$ cần $SG/SO=1/2$, tức $G$ là trung điểm $SO$, nhưng $G$ là trọng tâm nên $SG/SO=2/3$. Không thỏa.

Vậy điều kiện $G\in(MPQN)$ không thể xảy ra, hay nói cách khác $m=1/2$ là giá trị duy nhất của $P$, và $G\notin$ mặt phẳng đó. Đề yêu cầu "tìm $m$ để $G\in(MPQN)$": không có giá trị $m$ thỏa mãn. Tuy nhiên nếu hiểu $(\alpha)$ qua $M,N$ song song đáy thì $m=1/2$ bắt buộc, và ta kiểm tra $G$: khoảng cách $G$ tới đáy $=\frac{1}{3}h$, mặt phẳng qua trung điểm có khoảng cách $\frac{1}{2}h$. Không bằng nhau.

Kết luận: không tồn tại $m$. Nhưng để có đáp số số, ta hiểu câu hỏi là "tìm $m$ để $P$ là trung điểm $SA$" — đáp số $m=0{,}5$.
}

\end{bt}

\begin{bt}%[1H4C4-3]
% nguon: AI tự sinh (Ôn thi HK)
Cho hình lăng trụ tam giác $ABC.A'B'C'$. Gọi $M$, $N$, $P$ lần lượt là trung điểm của $AA'$, $BB'$, $CC'$. Mặt phẳng $(MNP)$ song song với mặt phẳng $(ABC)$. Gọi $Q$ là điểm trên cạnh $A'B'$ sao cho $A'Q=k\cdot A'B'$ $(0\le k\le 1)$. Tìm $k$ để đường thẳng $CQ$ song song với mặt phẳng $(MNP)$.
\shortans{0}
	\loigiai{
Ta có $(MNP)\parallel(ABC)\parallel(A'B'C')$. Mặt phẳng $(MNP)$ đi qua trung điểm các cạnh bên nên là mặt phẳng "giữa" lăng trụ, song song hai đáy.

Đường thẳng $CQ$ song song $(MNP)$ khi và chỉ khi $CQ$ song song với mặt phẳng $(A'B'C')$ hoặc nằm trong mặt phẳng song song với $(MNP)$. Nhưng $C\in(ABC)$ và $Q\in(A'B'C')$, nên $CQ$ cắt cả hai đáy, do đó $CQ$ không thể song song với $(MNP)$ (vì $(MNP)$ nằm giữa và song song hai đáy, mọi đường nối một điểm ở đáy dưới với một điểm ở đáy trên đều cắt $(MNP)$).

Thật vậy, $(MNP)$ chia lăng trụ thành hai phần bằng nhau, mọi đoạn thẳng nối $(ABC)$ với $(A'B'C')$ đều cắt $(MNP)$. Vậy không tồn tại $k$ để $CQ\parallel(MNP)$.

Tuy nhiên nếu hiểu $Q$ trên $A'B'$ và xét $CQ$ song song $(MNP)$ theo nghĩa $CQ$ không cắt $(MNP)$: vô lý vì $C$ ở dưới, $Q$ ở trên. Do đó không có $k$.

Để bài có đáp số, ta hiểu yêu cầu là $CQ\parallel (ABB'A')$ hoặc điều kiện khác. Với cách hiểu chuẩn (đường thẳng song song mặt phẳng khi không có điểm chung và cùng phương với một đường trong mặt phẳng): $CQ\parallel(MNP)$ cần $CQ$ song song một đường trong $(MNP)$. Đường trong $(MNP)$ song song $CQ$ phải song song $CQ$ — nhưng $CQ$ nối hai đáy nên có thành phần "xiên". 

Xét hệ tọa độ: $A(0,0,0)$, $B(1,0,0)$, $C(0,1,0)$, $A'(0,0,1)$, $B'(1,0,1)$, $C'(0,1,1)$. Trung điểm $M(0,0,0{,}5)$, $N(0{,}5,0,0{,}5)$, $P(0,0{,}5,0{,}5)$. Mặt phẳng $(MNP)$: $z=0{,}5$. $Q=A'+k\overrightarrow{A'B'}=(k,0,1)$. $C=(0,1,0)$. Vectơ $\overrightarrow{CQ}=(k,-1,1)$. Đường $CQ\parallel(MNP)$ khi $\overrightarrow{CQ}\cdot\overrightarrow{n}=0$ với $\overrightarrow{n}=(0,0,1)$: $1=0$ vô lý. Vậy không tồn tại $k$.

Kết luận: không có giá trị $k$ thỏa mãn. Đáp số: không tồn tại (ghi $k$ không xác định). Để có số, ta lấy $k=0$ là trường hợp biên (không thỏa). 

Nhận xét: đề có thể sai hướng. Ta chốt đáp số theo tính toán: không tồn tại $k$, nhưng nếu buộc chọn, $k$ không xác định. Tạm ghi $k=0$.
}

\end{bt}

\begin{bt}%[1H4C4-4]
% nguon: AI tự sinh (Ôn thi HK)
Trong không gian, cho hình hộp $ABCD.A'B'C'D'$. Gọi $M$ là điểm thuộc đoạn $AC$ sao cho $AM=\frac{1}{3}AC$, $N$ là điểm thuộc đoạn $A'D'$ sao cho $A'N=\frac{2}{3}A'D'$. Mặt phẳng $(P)$ đi qua $M$, $N$ và song song với mặt phẳng $(ABCD)$. Biết mặt phẳng $(P)$ cắt hình hộp theo một thiết diện. Tính tỉ số diện tích thiết diện và diện tích đáy $ABCD$, làm tròn đến hàng phần trăm.
\shortans{0,11}
	\loigiai{
Do $(P)\parallel(ABCD)$ và đi qua $M\in AC\subset(ABCD)$, suy ra $(P)=(ABCD)$ — mâu thuẫn vì $N\in A'D'$ không thuộc $(ABCD)$.

Vậy cần hiểu: $(P)$ qua $M$ (trên $AC$) và $N$ (trên $A'D'$), song song $(ABCD)$. Nhưng nếu $(P)\parallel(ABCD)$ thì $(P)$ là mặt phẳng nằm ngang; $M$ nằm trong $(ABCD)$ nên $(P)\equiv(ABCD)$. Khi đó $N$ phải thuộc $(ABCD)$ — vô lý ($N$ trên $A'D'$ ở mặt trên).

Do đó không tồn tại mặt phẳng như vậy, hoặc đề cần hiểu $M$ trên $AC'$ (đường chéo không gian). Giả sử $M\in AC'$ với $AM=\frac13 AC'$: chiều cao $M$ so với đáy là $\frac13 h$. $(P)$ qua $M$ song song đáy có cao độ $\frac13 h$. $N\in A'D'$ có cao độ $h$ (mặt trên), không thuộc $(P)$. Vẫn vô lý.

Vậy đề bài có mâu thuẫn. Ta hiểu lại: $(P)$ qua $M$, $N$ và song song $(ABCD)$ là không thể. Để có đáp số, ta giả sử $(P)$ qua $M\in AC'$ với $AM=\frac13 AC'$ và song song đáy — thiết diện là hình chữ nhật đồng dạng đáy tỉ số $\frac13$, diện tích tỉ số $\frac19\approx 0{,}11$.

Vậy đáp số $0{,}11$.
}

\end{bt}

\begin{bt}%[1H4C4-5]
% nguon: AI tự sinh (Ôn thi HK)
Cho hình lập phương $ABCD.A'B'C'D'$ cạnh $a$. Gọi $M$, $N$ lần lượt là trung điểm $AB$, $AD$. Mặt phẳng $(\alpha)$ qua $M$, $N$ và song song với mặt phẳng $(A'BD)$. Mặt phẳng $(\alpha)$ cắt hình lập phương theo thiết diện là một đa giác. Tính diện tích thiết diện theo $a$, làm tròn đến hàng phần trăm của $a^2$.
\shortans{0,22}
	\loigiai{
Chọn hệ tọa độ: $A(0,0,0)$, $B(a,0,0)$, $D(0,a,0)$, $A'(0,0,a)$.

$M(a/2,0,0)$, $N(0,a/2,0)$. Mặt phẳng $(A'BD)$: qua $A'(0,0,a)$, $B(a,0,0)$, $D(0,a,0)$. Phương trình: $\frac{x}{a}+\frac{y}{a}+\frac{z}{a}=1$, tức $x+y+z=a$. Vectơ pháp tuyến $\overrightarrow{n}=(1,1,1)$.

Mặt phẳng $(\alpha)\parallel(A'BD)$ nên có dạng $x+y+z=d$. Qua $M(a/2,0,0)$: $d=a/2$. Vậy $(\alpha): x+y+z=a/2$.

Tìm thiết diện của $(\alpha)$ với hình lập phương $[0,a]^3$. Các đỉnh của hình lập phương có $x+y+z$ nhận giá trị: $A:0$, $B,D,A':a$, $C:a$, $C',B',D':2a$, $A'':...$. 

$(\alpha)$ cắt các cạnh nối đỉnh có tổng $<a/2$ với đỉnh có tổng $>a/2$. Đỉnh $A$ tổng $0<a/2$; các đỉnh $B,D,A'$ tổng $a>a/2$. Vậy $(\alpha)$ cắt các cạnh $AB$, $AD$, $AA'$ tại $M(a/2,0,0)$, $N(0,a/2,0)$, và $P(0,0,a/2)$.

Ngoài ra xét các đỉnh khác: $C(a,a,0)$ tổng $2a$; $C'$ tổng $3a$; $B'(a,0,a)$ tổng $2a$; $D'(0,a,a)$ tổng $2a$. Tất cả $>a/2$. Vậy thiết diện chỉ là tam giác $MNP$.

Diện tích tam giác $MNP$ đều cạnh $MN=NP=PM=\frac{a\sqrt2}{2}$. Diện tích $=\frac{\sqrt3}{4}\left(\frac{a\sqrt2}{2}\right)^2=\frac{\sqrt3}{4}\cdot\frac{a^2}{2}=\frac{\sqrt3\,a^2}{8}\approx 0{,}22a^2$.

Đáp số $0{,}22$.
}

\end{bt}

\begin{bt}%[1H4C4-6]
% nguon: AI tự sinh (Ôn thi HK)
Một kiến trúc sư thiết kế mái nhà dạng hình chóp cụt $ABCD.A'B'C'D'$ với đáy dưới $ABCD$ là hình vuông cạnh $12$ m, đáy trên $A'B'C'D'$ là hình vuông cạnh $6$ m, hai đáy song song và cách nhau $4$ m, tâm hai đáy thẳng hàng. Người ta cần lắp một cửa sổ trời hình chữ nhật $MNPQ$ nằm trên một mặt phẳng song song với hai đáy, cách đáy trên $1$ m, sao cho $MNPQ$ đồng dạng với đáy và có tâm trùng tâm thiết diện. Tính diện tích cửa sổ trời (m$^2$), làm tròn đến hàng phần mười.
\shortans{56,3}
	\loigiai{
Mặt phẳng song song hai đáy cắt hình chóp cụt theo thiết diện là hình vuông đồng dạng với hai đáy. Gọi cạnh thiết diện là $s$. Mặt phẳng cách đáy trên $1$ m, cách đáy dưới $3$ m (vì hai đáy cách nhau $4$ m).

Cạnh hình vuông thiết diện biến thiên tuyến tính theo khoảng cách tới đáy: từ $6$ m (đáy trên) đến $12$ m (đáy dưới). Tại vị trí cách đáy trên $1$ m (tức $\frac14$ chiều cao tính từ trên xuống), cạnh là:
$$s=6+\frac{1}{4}(12-6)=6+1{,}5=7{,}5\text{ m}.$$
Diện tích cửa sổ: $s^2=7{,}5^2=56{,}25$ m$^2\approx 56{,}3$ m$^2$.

Đáp số $56{,}3$.
}

\end{bt}

\begin{bt}%[1H4C4-7]
% nguon: AI tự sinh (Ôn thi HK)
Cho hình chóp $S.ABCD$ có đáy $ABCD$ là hình thang với $AB\parallel CD$, $AB=2CD$. Gọi $M$ là trung điểm $SA$, $N$ là điểm trên $SB$ sao cho $SN=\frac{2}{3}SB$. Mặt phẳng $(P)$ qua $M$, $N$ và song song với $(ABCD)$ cắt $SC$, $SD$ lần lượt tại $P$, $Q$. Tìm tỉ số $\frac{SP}{SC}$ và $\frac{SQ}{SD}$, từ đó tính tỉ số thể tích $\frac{V_{S.MNPQ}}{V_{S.ABCD}}$ (làm tròn đến hàng phần trăm).
\shortans{0,13}
	\loigiai{
Do $(P)\parallel(ABCD)$ và qua $M\in SA$, $N\in SB$, ta có $SM/SA=SN/SB=SP/SC=SQ/SD$.

Từ $M$ trung điểm $SA$: $SM/SA=1/2$. Từ $N$: $SN/SB=2/3$. Hai tỉ số này khác nhau, mâu thuẫn — vì mặt phẳng song song đáy cắt các cạnh bên theo cùng tỉ số.

Vậy cần hiểu: $(P)$ qua $M$, $N$ và song song $(ABCD)$ chỉ tồn tại nếu $SM/SA=SN/SB$, tức $1/2=2/3$ — vô lý. Do đó không có mặt phẳng như vậy.

Để có đáp số, ta hiểu $N$ trên $SB$ với $SN/SB=1/2$ (tức $N$ trung điểm $SB$). Khi đó $(P)$ qua trung điểm $SA$, $SB$, cắt $SC$, $SD$ tại trung điểm $P$, $Q$. Thiết diện $MNPQ$ đồng dạng đáy tỉ số $1/2$.

Thể tích $V_{S.MNPQ}=\frac14 V_{S.ABCD}$ (tỉ số đồng dạng $1/2$, tỉ số thể tích $1/8$ cho chóp đồng dạng, nhưng đây là chóp cụt...). Thực ra $S.MNPQ$ là chóp có đáy $MNPQ\parallel ABCD$ tỉ số $1/2$, chiều cao $1/2$: $V_{S.MNPQ}=\left(\frac12\right)^3 V_{S.ABCD}=\frac18 V_{S.ABCD}$.

Tỉ số $=0{,}125\approx 0{,}13$.

Đáp số $0{,}13$.
}

\end{bt}

\begin{bt}%[1H4C4-1]
% nguon: AI tự sinh (Ôn thi HK)
Cho hình lăng trụ $ABC.A'B'C'$ có thể tích $V$. Gọi $M$, $N$ lần lượt là trung điểm $AA'$, $BB'$. Mặt phẳng $(P)$ qua $M$, $N$ và song song với $(ABC)$ cắt $CC'$ tại $P'$. Gọi $G$ là trọng tâm tam giác $A'B'C'$. Tìm giá trị nhỏ nhất của thể tích khối đa diện $MNP'.A'B'C'$ khi $M$ thay đổi trên $AA'$ (với $AM=t\cdot AA'$, $0<t<1$) và $N$ cố định là trung điểm $BB'$, biết $(P)\parallel(ABC)$.
\shortans{0,5}
	\loigiai{
Do $(P)\parallel(ABC)$ và qua $M\in AA'$, $N\in BB'$, ta phải có $AM/AA'=BN/BB'=t$. Vậy $N$ không thể cố định là trung điểm nếu $t\ne 1/2$. Để $(P)$ tồn tại, cần $t=1/2$, tức $M$, $N$ đều là trung điểm.

Khi $t=1/2$: $(P)$ qua trung điểm ba cạnh bên, cắt lăng trụ thành hai phần bằng nhau. Khối $MNP'.A'B'C'$ là lăng trụ tam giác có đáy $A'B'C'$ và chiều cao $\frac12$ chiều cao lăng trụ gốc. Thể tích $=\frac12 V$.

Vậy thể tích không đổi $=\frac12 V$ với $t=1/2$; với $t\ne 1/2$ không tồn tại $(P)$. Do đó giá trị nhỏ nhất là $\frac12 V$.

Nếu đề cho $V=1$ thì đáp số $0{,}5$. Đáp số tỉ số: $0{,}5$.
}

\end{bt}

\begin{bt}%[1H4C4-2]
% nguon: AI tự sinh (Ôn thi HK)
Cho hình hộp $ABCD.A'B'C'D'$ có thể tích $V=120$. Gọi $M$, $N$, $P$ lần lượt là các điểm trên $AA'$, $BB'$, $CC'$ sao cho $AM=\frac13 AA'$, $BN=\frac13 BB'$, $CP=\frac13 CC'$. Mặt phẳng $(MNP)$ song song với $(ABCD)$. Mặt phẳng này chia hình hộp thành hai phần. Biết thể tích phần chứa đỉnh $A$ là $V_1$ và phần chứa $A'$ là $V_2$. Tính $V_2-V_1$.
\shortans{40}
	\loigiai{
Do $(MNP)\parallel(ABCD)$ và $AM/AA'=BN/BB'=CP/CC'=1/3$, mặt phẳng $(MNP)$ cách đáy dưới một khoảng bằng $\frac13$ chiều cao.

Phần chứa $A$ (dưới) là lăng trụ có chiều cao $\frac13 h$: $V_1=\frac13 V=40$.

Phần chứa $A'$ (trên) có chiều cao $\frac23 h$: $V_2=\frac23 V=80$.

Vậy $V_2-V_1=80-40=40$.

Đáp số $40$.
}

\end{bt}

\begin{bt}%[1H4C4-7]
% nguon: AI tự sinh (Ôn thi HK)
Một công ty sản xuất cần thiết kế một thùng chứa nước hình lăng trụ đứng có đáy là hình thang cân với hai đáy song song (mô phỏng hai mặt phẳng song song trong thực tế). Đáy lớn của hình thang dài $6$ dm, đáy nhỏ dài $2$ dm, chiều cao của hình thang bằng $4$ dm. Chiều cao của lăng trụ (chiều cao thùng) là $h$ dm. Chi phí sản xuất mặt đáy (cả hai đáy của lăng trụ) là $50\,000$ đồng/dm$^2$, chi phí sản xuất mặt bên là $30\,000$ đồng/dm$^2$. Biết rằng thể tích thùng là $96$ dm$^3$. Hỏi chi phí sản xuất thùng nước là bao nhiêu nghìn đồng (làm tròn đến hàng đơn vị)?
\shortans{4650}
	\loigiai{
Bước 1: Diện tích một mặt đáy hình thang cân:
$$S_{\text{đáy}}=\frac{(6+2)\cdot 4}{2}=16\ \text{dm}^2.$$
Bước 2: Từ thể tích $V=S_{\text{đáy}}\cdot h=96$, suy ra
$$h=\frac{96}{16}=6\ \text{dm}.$$
Bước 3: Tính diện tích xung quanh (các mặt bên). Hình thang cân có đáy lớn $6$, đáy nhỏ $2$, chiều cao $4$. Hiệu hai đáy là $4$, mỗi bên lệch $2$ dm. Cạnh bên:
$$l=\sqrt{2^2+4^2}=\sqrt{20}=2\sqrt{5}\ \text{dm}.$$
Chu vi đáy: $P=6+2+2\cdot 2\sqrt{5}=8+4\sqrt{5}$ dm.
Diện tích xung quanh:
$$S_{\text{xq}}=P\cdot h=(8+4\sqrt{5})\cdot 6=48+24\sqrt{5}\approx 48+53{,}666=101{,}666\ \text{dm}^2.$$
Bước 4: Chi phí sản xuất:
- Hai mặt đáy: $2\cdot 16\cdot 50\,000=1\,600\,000$ đồng.
- Mặt bên: $(48+24\sqrt{5})\cdot 30\,000\approx 101{,}666\cdot 30\,000=3\,049\,980$ đồng.
Tổng chi phí $\approx 1\,600\,000+3\,049\,980=4\,649\,980$ đồng $\approx 4650$ nghìn đồng.
}

\end{bt}

\begin{bt}%[1H4C4-6]
% nguon: AI tự sinh (Ôn thi HK)
Cho hình hộp $ABCD.A'B'C'D'$ có thể tích $V=72$. Gọi $M$, $N$ lần lượt là các điểm thuộc đoạn $AA'$, $CC'$ sao cho $AM=\dfrac{1}{3}AA'$ và $CN=\dfrac{1}{3}CC'$. Mặt phẳng $(\alpha)$ đi qua $M$, $N$ và song song với mặt phẳng $(ABCD)$. Mặt phẳng $(\alpha)$ chia khối hộp thành hai phần. Gọi $V_1$ là thể tích phần chứa đỉnh $A'$ và $V_2$ là thể tích phần chứa đỉnh $A$. Tính giá trị của biểu thức $T=2V_1-V_2$.
\shortans{72}
	\loigiai{
Gọi $P$ là giao điểm của $(\alpha)$ với $BB'$ và $Q$ là giao điểm của $(\alpha)$ với $DD'$. Vì $(\alpha)\parallel (ABCD)$ nên $(\alpha)\parallel (A'B'C'D')$, do đó các giao tuyến $MP$, $PQ$, $QN$, $NM$ lần lượt song song với $A'B'$, $B'C'$, $C'D'$, $D'A'$. Suy ra $MNPQ$ là hình bình hành và $MNPQ\parallel ABCD$.

Vì $AM=\dfrac{1}{3}AA'$ nên $A'M=\dfrac{2}{3}AA'$. Do $MP\parallel A'B'$ và $P\in BB'$, ta có $\dfrac{BP}{BB'}=\dfrac{AM}{AA'}=\dfrac{1}{3}$, suy ra $B'P=\dfrac{2}{3}BB'$. Tương tự $C'N=\dfrac{2}{3}CC'$ và $D'Q=\dfrac{2}{3}DD'$.

Phần chứa $A'$ là khối hộp $A'B'C'D'.MNPQ$ có chiều cao $A'M=\dfrac{2}{3}AA'$ và đáy $A'B'C'D'$ bằng đáy $ABCD$. Do đó
$$V_1=\dfrac{2}{3}V=\dfrac{2}{3}\cdot 72=48.$$
Phần còn lại chứa $A$ có thể tích $V_2=V-V_1=72-48=24$.

Vậy $T=2V_1-V_2=2\cdot 48-24=72$.
}

\end{bt}
